\section{Sensors and Actuators}
\subsection{Measurements}
Sensors model a state of the environment through (time discrete) observations.\hfill
\begin{figure}[h]
    \centering
    \includegraphics[width=.6\textwidth]{measure.png}
\end{figure}

The (analogue) signal captured by sensors is then digitised, i.e. converted from being continuous in \f{t,x} to being dicontinuous in both dimensions. To achieve this, different conversion strategies exist, all differing in conversion speed, accuracy of the values, and amplitude resolution after quantisation.\hfill
\begin{figure}[h!]
    \centering
    \includegraphics[width=.7\textwidth]{signals.png}
\end{figure}
% \noindent
\subsection{Parallel/Direct Conversion}
\begin{wrapfigure}{r}{.3\textwidth}
    % \centering
    % \includegraphics[width=.98\linewidth]{parallel.png}
    \raisebox{0pt}[\dimexpr\height-0.8\baselineskip\relax]{\includegraphics[width=.98\linewidth]{parallel.png}}
\end{wrapfigure}
Parallel converters work by comparing input voltage to a reference voltage using \f{2^n-1} parallel comparators.

The example on the right shows a 3-Bit converter. This approach is typically limited to \f{2^{N=8}=256} states and achieves frequencies greater than 1GHz.

To convert an n-bit digital signal into analogue, one can then use a parallel DAC as pictured below (n switches). The problem is that digital switches could be asynchronous, causing undefined u\f{_{\textrm{out}}}.
\vfill
\begin{figure}[h]
    \centering
    \includegraphics[width=.6\textwidth]{dac.png}
\end{figure}
\vfill
\clearpage

\subsection{Successive Approximation}
\begin{wrapfigure}{r}{.3\textwidth}
    % \centering
    % \includegraphics[width=.98\linewidth]{successive.png}
    \raisebox{0pt}[\dimexpr\height-2\baselineskip\relax]{\includegraphics[width=.98\linewidth]{successive.png}}
\end{wrapfigure}
Successive Approximation is a sequential conversion technique that works according to the \b{weighting principle}. 

A DAC creates different reference voltages to be compared to the input. SA can resolve 24bits and upwards at a frequency of around 1MHz.

\subsection{Sample \& Hold}
The idea is to charge up to u\f{_{\textrm{in}}} as fast as possible (sample) and hold u\f{_{\textrm{out}}} constant for as long as possible.
\begin{figure}[h]
    \centering
    \includegraphics[width=.65\textwidth]{sh.png}
\end{figure}

\subsection{Dual/Multi-Slope Integrating Converter}
The main idea behind this technique is to convert and unknown voltage into a time delay, which is then encoded. It is accurate, produces low noise, has a linear transfer function, but is rather slow with a max. frequency of around 0.1-1kHz.
\begin{enumerate}
    \item Integrate u\f{_{\textrm{in}}} for fixed time (until counter flips to 0), results in u\f{_{\textrm{C}}}
    \item Apply u\f{_{\textrm{ref}}} with opposite polarity to integrator until 0V, thus comparator zeros and blocks clock at gate to maintain count result 
    \item Encode u\f{_{\textrm{in}}} as a function of u\f{_{\textrm{ref}}}, run-up time \f{t_0} and run-down time \f{t_U}
\end{enumerate}

\subsection{Aliasing}
Aliasing describes the effect when e.g. car wheels seem to spin backwards while accelerating. The same effect can also occure when sampling from an input signal.

To avoid aliasing, sampling intervals \f{T_s = 1/f_s} must be carefully chosen so that they can represent the original signal. A signal with maximum frequency of \f{f_{\textrm{in}}} must be sampled at \f{f_s\geq 2f_{\textrm{in}}}. Aliasing orrurs when \f{f_s < 2 f_{\textrm{in}} \leftrightarrow 0.5f_s < f_{\textrm{in}}}.

For a signal with two frequency components, one must sample at \f{f_s \geq 2(f_b - f_a)} to prevent aliasing.\hfill
\newline

More generally, given a sampling rate \f{f_s}, two frequencies \f{f} and \f{f'} are aliases of each other if \f{\exists k\in\mathbb{N}: f' = f + kf_s}. The \b{Nyquist Theorem} states, that the only range in which the original frequency can correctly be sampled/reconstructed is the \b{Nyquist bandwidth} from 0 to \f{0.5f_s}. 

Example: sampling at \f{0.2f_s \rightarrow f_s \geq 5f_{\textrm{in}}\rightarrow\textrm{ no aliasing}}.

\subsection{Aliasing Filters}
\b{Baseband sampling:} All signals of interest lie within Nyquist bandwidth (first Nyquist zone). Images at higher frequencies are filtered.\hfill\\[.5em]
\b{Undersampling:} Sampling signals outside first Nyquist zone. Image in first Nyquist zone contain all information in frequency-reversed order. Sampled signal may lie in any Nyquist zone. Image in first Nyquist zone retains all information. For this the Nyquist criterion is updated to \f{f_s\geq2\Delta f} where \f{\Delta f} is the bandwidth of the signal of interest.

\clearpage
\begin{figure}[t]
    \centering
    \includegraphics[width=.5\textwidth]{nyquist.png}
\end{figure}
\begin{figure}[t]
    \centering
    \includegraphics[width=.5\textwidth]{nyquist2.png}
\end{figure}





